% robustness, validation and literature comparison (all numbers from generated macros)
\section{How much can these results be trusted?}\label{sec:checks}

\paragraph{Robustness.} All checks were re-run under the final configuration; none of the mid-evaluation
counts is reused. We varied one choice at a time: normalisation (x/max vs min-max), exposure definition
(level, level + change, change only), sensitivity (literature, blended as at the mid-evaluation, empirical
only), rice $K_y$, nutrient weighting, heat thresholds ($\pm$1--3\,\si{\celsius}), the new crops' proxy
parameters and their calendars. Wheat is first in \robTopFfour\ of \robNFfour\ \Fd\ rankings (both
weightings) and in \robTopAll\ of \robN\ rankings over all four index forms. The exceptions are the
``empirical sensitivity only'' tests, which rest on regressions that detect a heat-damage effect in only \empHeatDetected\ of \empN\ pairs (\empHeatCountry\ with a robust country-specific estimate, the rest from pooled panels). With 2,000 random weightings (level vs change, heat vs water, nutrient weights; seed 20261005) wheat
ranks first in \mcTopShare\% of draws and \mcSecond\ in \mcSecondShare\%; these are shares of weightings,
not probabilities that a crop is truly most vulnerable (Figure~\ref{fig:robust}). What does \emph{not}
survive: the order of positions two to six, which moves with the weighting and with the pulse parameters.

\paragraph{Validation.} Against Zhao et al.'s per-degree losses \citep{zhao2017} the climate part of the
index has Spearman $\rho$ = \zhaoRhoClimate\ and the full index \zhaoRhoFfour\ (four crops only, so
descriptive). Across the \shockNpairs\ pairs, the climate part is positively associated with how often yield
fell more than one standard deviation below trend in 2001--2024 ($\rho$ = \shockRho, one-sided permutation
$p$ = \shockP). At the mid-evaluation, with 26 pairs, this test was null; the nine-crop result should be read
as an association in a larger and more varied sample, not as proof that the index predicts shocks. The
fitted year-to-year climate response correlates positively with observed 2001--24 yield anomalies in
\oosPos\ of \oosN\ pairs with enough data (median $r$ = \oosMedian).

\paragraph{Comparison with published studies.} Of \litN\ comparisons of signs and rank orders
(Table~\ref{tab:lit}), \litAgree\ agree, \litPartial\ partly agree and \litDiffer\ differ. The
disagreements are informative: \citet{lobell2011} attribute Russian wheat losses to temperature, while our
monthly threshold measure sees mainly rainfed water deficit there; and maize suitability studies
\citep{ramirez2017} project losses in Africa where our single climate model projects a wetter South African
growing season. Agreement is supporting context, not proof of accuracy.

\begin{table}[p]
\centering\scriptsize
\caption{Our results next to published findings.}\label{tab:lit}
\begin{minipage}{\textwidth}\centering
\adjustbox{max width=\textwidth}{\begin{tabular}{p{0.9cm}p{3.0cm}p{4.4cm}p{4.6cm}p{2.0cm}}
\toprule
 & Study & Published finding & Our comparable result & Verdict \\
\midrule
Z1 & Zhao et al. 2017 PNAS & yield change per +1 C: maize -7.4, wheat -6.0, rice -3.2, soybean -3.1 (n.s.) -$>$ maize $>$ wheat $>$ rice $>$ soybean & climate part E x S (equal weights): maize $>$ wheat $>$ soybean $>$ rice; Spearman rho with Zhao = +0.8 (n = 4) & agree (order of maize, wheat) \\
Z2 & Zhao et al. 2017 PNAS (Fig. 3) & wheat loss per +1 C: India -9.1, Russia -7.8, USA -5.5, China -2.6 \% & wheat heat exposure E\_heat: India 0.115, Russia 0.015, USA 0.015, China 0.002 & agree (heat) \\
L1 & Lobell, Schlenker \& Costa-Roberts 2011 Science & Russia wheat production about -15 \% from temperature trends & wheat-Russia F4 = 0.1315 (top pair); heat part 0.0027 vs water part 0.1288 & differ (mechanism) \\
R1 & Ray et al. 2015 Nat. Commun.; Ray et al. 2019 PLoS ONE & climate explains 32-39 \% of yield variability; rice least & rice climate part E x S = 0.0247, lowest of the four (wheat 0.1077, maize 0.1345, soybean 0.0830) & not directly comparable (different outcome) \\
D1 & Deryng et al. 2014 ERL (Table 1) & maize -10.5, wheat -10.1, soybean -6.9 \% (RCP4.5 with CO2: -4.9, +13.0, +10.8) & our conditional yield change (mean over countries, no CO2): maize -2.5 \%, wheat -3.5 \%, soybean +0.7 \% & sign agrees for maize; sign agrees for wheat; sign differs for soybean; our magnitudes smaller \\
C1 & Challinor et al. 2014 Nat. Clim. Change & $>$ 70 \% of projections show yield losses & 76 \% of our 59 pairs show a climate-driven yield loss (literature accounting); cereals only: 74 \% & agree \\
RC & Ramirez-Cabral et al. 2017 Sci. Rep. & suitable area: Africa -7 to -11 \%, South America -5 \%, India marginal areas become unsuitable & maize change in F4 2041-60 vs 1971-2000: South Africa -0.0127, Brazil +0.0019, India -0.0002 & differ \\
RP/HU & Rippke et al. 2016 Nat. Clim. Change; Hummel et al. 2018 Sci. Rep. & beans -10 \% area needing transformation in Uganda/Tanzania by the 2050s (RP); most SE African bean areas unsuitable by 2040-69, drought-driven (HU) & bean-Tanzania F4 change +0.0188; conditional yield change -17.7 \% & agree (direction) \\
Q1 & Qian et al. 2019 ERL (Table 2) & -0.8 to -5.7 \% production at +2 C with CO2 fixed (+3 to +11 \% with CO2) & rapeseed-Canada conditional yield change -3.2 \% (no CO2); flowering T\_eff change +3.1 C & agree (sign) \\
K1 & Kalra et al. 2008 Curr. Sci. & yields fall with warmer seasons: chickpea -0.53 to -3.01, mustard -0.92 to -2.01 q/ha per C & India 2041-60: chickpea T\_eff +1.3 C, yield -0.2 \%; rapeseed-mustard T\_eff +1.6 C, yield -2.2 \% & agree (sign) \\
BE/NE & Beach et al. 2019 Lancet Planet. Health; Nelson et al. 2018 Nat. Sustain. & South Asia iron availability below the recommended intake; micronutrients, not calories, are the binding constraint & India 2021-23 iron supply / requirement = 0.89 (15 \% bioavailability: 1.00); energy 1.00, protein 1.00 & agree (main requirement) \\
M1 & Myers et al. 2014 Nature & iron concentration: wheat -5.1, rice -5.2, soybean -4.1 \% (significant) & applied to India's 2021-23 iron supply shares, these declines alone would lower national iron supply by 2.2 \% (an additional pathway the index excludes) & additional pathway \\
\bottomrule
\end{tabular}
}
\par\vspace{2pt}{\footnotesize Signs and rank orders only; published values verified in docs/literature\_comparison\_sources.md; ours computed by code/22\_literature\_comparison.py.}
\end{minipage}

\end{table}
